\section{Introduction and objectives}

Shor's algorithm factors integers in time polynomial in the binary length of the input,
with implications for the security of cryptographic systems based on the difficulty of
factorization \cite{en-shor1994,en-nist2024pqc,en-gidney2021}. Realizing this theoretical
advantage, however, requires error control compatible with circuit depth. Gate errors,
decoherence, and imperfect readout can alter the measurement distribution and compromise
the information needed to find the multiplicative order \cite{en-preskill2018}.

The thesis develops a simulation study aimed at establishing which interventions actually
improve the reliability of the computation, and relative to which baselines. The starting
point is the measurement of Shor's sensitivity to gate errors; three strategies follow:
outcome selection through post-processing, with an ablation study of the
machine-learning contribution; quantum error correction, with surface codes and qLDPC
codes; and the reduction of the number of gates through the approximate quantum Fourier
transform (AQFT).

The areas of intervention---outcome selection, information protection, and circuit cost
reduction---are linked by a common criterion: compare each solution with an explicit
baseline, isolate the contribution of individual components, and measure the benefit with
task-relevant metrics. The experimental campaigns were conducted separately; in
particular, the error-correction methods were not applied directly to the full Shor
circuit in a single simulation.
